% Part: sets-functions-relations
% Chapter: arithmetization
% Section: reflections

\documentclass[../../../include/open-logic-section]{subfiles}

\begin{document}

\olfileid{sfr}{arith}{ref}
\olsection{کجھ فلسفیانہ غور}

فنی تفصیلاں لئی اینا کافی اے۔ پر اوہناں نے حاصل کی کیتا؟

بھلا، لگ بھگ بے اختلاف طور تے اوہناں نے سانوں کجھ {سوہنی} خالص ریاضی
دِتی۔ نال ای کجھ ڈونگھیاں تصوری کامیابیاں وی سن۔ ایہ ویکھنا اک گہری
بصیرت سی کہ تکمیلیت دی خاصیت حقیقی تے ناطق عدداں وچکار فیصلہ کن فرق
بیان کردی اے۔ نال ای ڈیڈیکائنڈ کٹاں دے طور تے حقیقی عدداں دی صاف
بناوٹ، تجزیے دے موضوع نوں پکی بنیاد دیندی اے۔ اسی جان دے آں کہ
\emph{مکمل ترتیب دار میدان} دا تصور مربوط اے، کیوں جو کٹ عین ایسا
میدان بناؤندے نیں۔

ایس سب دے باوجود، سانوں ایہناں کامیابیاں بارے کجھ تحفظات سامنے لیانے چاہیدے نیں۔

پہلاں، ایہ صاف نہیں کہ حقیقی عدداں نوں کٹاں دے حوالے نال سوچنا، اوہناں
دیاں واقف (ممکنہ طور تے لامتناہی) اعشاری توسیعاں دے حوالے نال سوچن توں
کوئی \emph{ودھ} کڑا اے۔ حقیقی عدداں دی ایہ پچھلی ”بناوٹ“، کوشی لڑی
راہی حقیقی عدداں دی بناوٹ نال کجھ رل دی اے؛ پر اصل وچ ایہ سترھویں صدی
دے مُڈھ توں اگے ریاضی داناں نوں بنیادی طور تے معلوم سی (ویکھو
\olref[cauchy]{sec})۔ کڑائی وچ اصل اضافہ ایس ادراک توں آیا کہ حقیقی
عدداں کول تکمیلیت دی خاصیت اے؛ حقیقی عدداں نوں خاص سیٹاں دے طور تے
بنا سکنا شاید آپنے آپ وچ اینا دلچسپ نہیں۔

ایہ گل ہور وی گھٹ صاف اے کہ صحیح عدداں یا ناطق عدداں دی (بہت سوکھی)
اعدادی تشکیل، اوہناں شعبیاں وچ کڑائی ودھاؤندی اے۔ ایتھے اک سادہ گل
کرنی فائدہ مند اے۔ صحیح عدداں نوں قدرتی عدداں دیاں ترتیب وار جوڑیاں دے
ہم ارزیت طبقیاں دے طور تے \emph{بنا} کے، فیر ناطق عدداں نوں صحیح عدداں
دیاں ترتیب وار جوڑیاں دے ہم ارزیت طبقیاں دے طور تے بنا کے، تے فیر
حقیقی عدداں نوں ناطق عدداں دے سیٹاں دے طور تے بنا کے، اسی فوراً ایہناں
بناوٹاں نوں \emph{بھل جاندے} آں۔ خاص طور تے: کوئی وی ریاضیاتی ثبوت
وچ ایہناں بناوٹاں نوں \emph{ورتنا} نہیں چاہے گا (بے شک اوس ثبوت توں
سوا جیہڑا وکھاوے کہ بناوٹاں نے اوہو جیہا ورتاؤ کیتا جیہو جیہا اوہناں
نوں کرنا چاہیدا سی)۔ سیٹاں دے سیٹاں دے سیٹاں دے سیٹاں دے سیٹاں دے
سیٹاں دے سیٹاں دے قدرتی عدداں بارے گل کرن نالوں، سیدھا حقیقی عدد بارے
گل کرنا بہت سوکھا اے۔
%(کیوں جو حقیقی عدد ناطق عدداں دے سیٹاں دے طور تے بنائے گئے سن؛ تے ایہ صحیح عدداں دیاں ترتیب وار جوڑیاں دے ہم ارزیت طبقیاں، یعنی صحیح عدداں دے سیٹاں دے سیٹاں دے سیٹاں، دے طور تے بنائے گئے سن؛ وغیرہ۔)
%2/3 = [2,3]
%	یعنی {<2,3>, ... }
%	{ {{2},{2,3}}, ... }
%
%حقیقی عدد
%	ناطق عدداں دے سیٹ نیں
%
%ناطق عدد
%	صحیح عدداں دے سیٹاں دے سیٹاں دے سیٹ نیں
%
%صحیح عدد
%	قدرتی عدداں دے سیٹاں دے سیٹاں دے سیٹ نیں

سب توں ودھ شک ایس گل اُتے اے کہ ایہ تعریفاں سانوں دسدیاں نیں کہ صحیح
عدد، ناطق عدد یا حقیقی عدد \emph{مابعدالطبیعیاتی طور تے} \emph{کی نیں}۔
یعنی ایس گل اُتے شک اے کہ حقیقی عدد (مثال لئی) کجھ خاص سیٹ (سیٹاں دے
سیٹاں دے سیٹ\ldots) \emph{ای نیں}۔ ایسے نظریے دے راہ وچ وڈی رکاوٹ
ایہ اے کہ بناوٹ کئی وکھرے طریقیاں نال کیتی جا سکدی سی۔ حقیقی عدداں دے
معاملے وچ کجھ واقعی دلچسپ طور تے وکھریاں بناوٹاں نیں (ویکھو
\olref[cauchy]{sec})۔ پر کجھ وکھریاں بناوٹاں حاصل کرن دا بالکل معمولی
طریقہ ایہ اے: \olref[sfr][rel][ref]{sec} وانگ، اسی ترتیب وار جوڑیاں نوں
ذرا وکھرے ڈھنگ نال تعریف کر سکدے ساں؛ جے اسی ترتیب وار جوڑے دا ایہ
متبادل تصور ورتدے، تاں ساڈیاں بناوٹاں عین اوہنی کامیابی نال کم کردیاں،
پر آخر وچ سانوں وکھریاں شےواں ملدیاں۔ ایس لئی صحیح، ناطق تے حقیقی
عدداں دیاں کئی مقابل سیٹ-تھیوری بناوٹاں نیں۔ تے ہُن ایہ دعویٰ کرنا محض
من مانا (تے شرمندگی والا) ہووے گا کہ صحیح عدد (مثال لئی) \emph{ایہ}
سیٹ نیں، \emph{اوہ} نہیں۔ (جیویں \olref[sfr][rel][ref]{sec} وچ، ایہ
\citealt{Benacerraf1965} نال مشہور ہوئی دلیل دی اک مثال اے۔)

اک ہور نکتہ وی چکنا فائدہ مند اے: ساڈیاں بناوٹاں بارے کجھ بڑا
\emph{انجوکھا} اے۔ اسی قدرتی عدداں توں شروع کیتا۔ فیر اسی صحیح عدد
بناؤندے آں، تے ”صحیح عدداں دا $0$“ بناؤندے آں، یعنی
$ \equivrep{0,0}{\Intequiv}$۔ پر $0 \neq
\equivrep{0,0}{\Intequiv}$۔ بلکہ ساڈیاں بناوٹاں دے حساب نال کوئی وی
قدرتی عدد صحیح عدد \emph{نہیں}۔ پر ایہ بہت خلافِ وجدان لگدا اے۔ بلکہ
\olref[sfr][set][imp]{sec} وچ اسی بہت دلیل دِتے بغیر دعویٰ کیتا سی کہ
$\Nat \subseteq \Rat$۔ جے بناوٹاں سانوں عین \emph{دسدیاں} نیں کہ عدد
کی نیں، تاں ایہ دعویٰ صاف طور تے جھوٹا سی۔

سو کجھ پچھے ہٹ کے ویکھئیے، تاں اسی کتھے پہنچدے آں؟ سادی سیٹ تھیوری وچ
کم کردیاں تے قدرتی عدداں نوں پہلے توں ملیا من کے، اسی صحیح، ناطق تے
حقیقی عدداں نوں کجھ خاص سیٹاں دے طور تے \emph{ورت} سکدے آں۔ ایس معنٰی
وچ اسی ایہناں ہستیاں دے نظرییاں نوں سیٹ تھیوری دے اندر \emph{سما}
سکدے آں۔ پر ایس سماوٹ دی فلسفیانہ اہمیت بالکل سیدھی سادی نہیں۔

بے شک ایہناں وچوں کوئی گل آخری فیصلہ نہیں!{} نکتہ صرف ایہ اے۔ ایہ
وکھاؤن لئی کہ حقیقی عدداں دی اعدادی تشکیل دی \emph{واقعی} گہری فلسفیانہ
اہمیت اے، کوئی اضافی \emph{فلسفیانہ} دلیل دینا ضروری ہووے گا۔

%ہور ایہ کہ: ایس پورے باب دوران اسی قدرتی عدداں نوں سادہ طور تے پہلے توں
%\emph{ملیا} منیا اے۔ پر اسی اوہناں نال کی کر سکدے آں؟ ایہ اگلے باب دا
%موضوع اے۔

\end{document}
